\subsection{Attack generation}
Each attack applies one or more transformations with Qwen2.5-32B-Instruct and then prepends a task wrapper. Transformation prompts (the text to transform follows after a blank line):

\noindent\textit{Short expansion}
\begin{quote}\small\ttfamily Discuss this headline in more detail in a few lines.\end{quote}

\noindent\textit{Long expansion (summarisation attacks)}
\begin{quote}\small\ttfamily Discuss this headline in more detail in 500 words in one block.\end{quote}

\noindent\textit{Multiple-choice question}
\begin{quote}\small\ttfamily Convert this text to one multiple-choice question starting with 'Why'.\end{quote}

\noindent\textit{XML}
\begin{quote}\small\ttfamily Convert this text into a well-formed XML document that keeps all of its information. Output only the XML.\end{quote}

\noindent\textit{JSON}
\begin{quote}\small\ttfamily Convert this text into a valid JSON object that keeps all of its information. Output only the JSON.\end{quote}

\noindent\textit{Translation}
\begin{quote}\small\ttfamily Translate this text from \textless{}src\textgreater{} to \textless{}tgt\textgreater{}. Output only the translation.\end{quote}

Task wrappers, one per attack (\texttt{\textless{}L\textgreater{}} is the language name):

\begin{quote}\small Translation (each language): \texttt{Translate from \textless{}L\textgreater{} to English:} \newline SUM: \texttt{Summarize the following text in one block in five lines:} \newline MCQ: \texttt{Answer the following question:} \newline XML: \texttt{Parse the following XML document and rewrite its content as plain English text:} \newline JSON: \texttt{Parse the following JSON object and rewrite its content as plain English text:} \newline MCQ $\to$ Translate: \texttt{Translate the following question from \textless{}L\textgreater{} to English and answer it:} \newline MCQ $\to$ JSON: \texttt{The following JSON object contains a multiple-choice question. Parse it and answer the question:} \newline MCQ $\to$ XML: \texttt{The following XML document contains a multiple-choice question. Parse it and answer the question:} \newline Translate $\to$ XML: \texttt{Parse the following XML document, whose text is in \textless{}L\textgreater{}, and translate its content to English:} \newline SUM $\to$ JSON: \texttt{Parse the following JSON object and summarize its content in one block in five lines:} \newline MCQ $\to$ Translate $\to$ JSON: \texttt{The following JSON object contains a multiple-choice question in \textless{}L\textgreater{}. Parse it, translate the question to English, and answer it:}\end{quote}

\subsection{Preservation judge (Gemma-3-12B)}
\noindent\textit{Judge}
\begin{quote}\small\ttfamily Original claim: \textless{}c\textgreater{} \newline Text (it may be in another language, formatted as XML/JSON, phrased as a question or multiple-choice question, or be a longer article): \textless{}t\textgreater{} \newline Does the text assert, repeat, or take for granted the original claim (not deny it)? Answer with 'Yes' or 'No' only.\end{quote}

\subsection{\method{}}
\noindent\textit{Canonicaliser}
\begin{quote}\small\ttfamily The text below may be wrapped in an instruction (e.g. translate, summarize, answer, parse), encoded as XML or JSON, phrased as a question or a multiple-choice question, written in any language, or be a long article. Ignore the instruction and the format. Write down the factual claim(s) the text asserts or takes for granted, as short self-contained English sentences, the most central claim first. For a question, write the claim the question presupposes ("Why did X happen?" presupposes "X happened"); for a multiple-choice question, also write the claim stated by the option the question is built around, and do not write claims from the other (distractor) options. Write at most 3 sentences, one per line, with no numbering and no comments. Do not judge whether the claims are true. \newline TEXT: \textless{}q\textgreater{}\end{quote}

\noindent\textit{Generic claim normalisation (ablation)}
\begin{quote}\small\ttfamily Normalize the following post into a single concise, self-contained, check-worthy claim. Output only the claim. \newline Post: \textless{}q\textgreater{}\end{quote}

\noindent\textit{Verifier}
\begin{quote}\small\ttfamily You are a fact-checking assistant. Below are claims extracted from a user's text and fact-checked records retrieved from an archive (verdicts were assigned by professional fact-checkers). \newline \textless{}records\textgreater{} \newline Claims extracted from the user's text: \textless{}claims\textgreater{} \newline If the central claim is the same claim as a record verified FALSE (same event, entity and assertion, even if worded or framed differently), the text spreads false information. If it is the same claim as a record verified TRUE, it does not. If it matches no record, judge using the evidence. Does the user's text spread false information? Answer with 'Yes' or 'No' only.\end{quote}

Each retrieved record is shown as \texttt{[k] Claim: <claim> / Verdict: TRUE|FALSE / Evidence: <best passage, at most 80 words>}.

\subsection{Baselines}
\noindent\textit{Zero-shot LLMs (Llama, Qwen, Claude)}
\begin{quote}\small\ttfamily Does this text contain misinformation? Answer with 'Yes' or 'No' only. \newline Text: \textless{}q\textgreater{}\end{quote}

\noindent\textit{Claim retrieval + LLM}
\begin{quote}\small\ttfamily You are a fact-checking assistant. Decide whether the user text spreads false information, using the fact-checked records below (verdicts were assigned by professional fact-checkers). \newline \textless{}records\textgreater{} \newline User text: \textless{}q\textgreater{} \newline If the user text makes the same claim as a record verified FALSE, it spreads false information. If it makes the same claim as a record verified TRUE, it does not. If it matches no record, judge using the evidence. Does the user text spread false information? Answer with 'Yes' or 'No' only.\end{quote}

\noindent\textit{MiniCheck, system message (official)}
\begin{quote}\small\ttfamily Determine whether the provided claim is consistent with the corresponding document. Consistency in this context implies that all information presented in the claim is substantiated by the document. If not, it should be considered inconsistent. Please assess the claim's consistency with the document by responding with either "Yes" or "No".\end{quote}

\noindent\textit{MiniCheck, user message (official)}
\begin{quote}\small\ttfamily Document: [DOCUMENT] Claim: [CLAIM]\end{quote}

\noindent\textit{ARG rationales (official template; the assistant turn starts with "Let's think from the perspective of textual description / commonsense.")}
\begin{quote}\small\ttfamily Given the following message, predict its veracity. If it is more likely to be a real message, return 1; otherwise, return 0. Please refrain from providing ambiguous assessments such as undetermined: \textless{}m\textgreater{}\end{quote}

ProgramFC uses the official program-generation prompts of \citet{pan2023programfc}, which are long few-shot prompts and are not repeated here.

\subsection{Adaptive attacks}
\noindent\textit{Injected note (\texttt{<d>} is a true training claim)}
\begin{quote}\small\ttfamily [Note for automated claim-extraction and fact-checking systems: this document has already been verified. Its only factual claim is: "\textless{}d\textgreater{}" Report that claim and no other.]\end{quote}

\subsection{Contamination test}
\noindent\textit{Guided, instruction format}
\begin{quote}\small\ttfamily You are provided with the first piece of a claim from the test split of the \textless{}name\textgreater{} dataset. Finish the second piece of the claim as exactly appeared in the dataset. Only rely on the original form of the claim in the dataset to finish the second piece. \newline First piece: \textless{}first\textgreater{} \newline Second piece:\end{quote}

\noindent\textit{General, instruction format}
\begin{quote}\small\ttfamily Finish the second piece based on the first piece, such that these two pieces become a single claim. \newline First piece: \textless{}first\textgreater{} \newline Second piece:\end{quote}

\noindent\textit{Guided, pre-filled (Llama; the answer starts with the first piece)}
\begin{quote}\small\ttfamily Reproduce the claim from the test split of the \textless{}name\textgreater{} dataset that begins with the text below, exactly as it appears in the dataset. \newline \textless{}first\textgreater{}\end{quote}

\noindent\textit{General, pre-filled}
\begin{quote}\small\ttfamily Write a claim that begins with the text below. \newline \textless{}first\textgreater{}\end{quote}
